\paragraph{Cone logarithmic coefficient at \(N_s=6\), \(n_D=2\) on the fixed-state replica (conditional; not closure).}
Keep the setting of the preceding paragraph and the fixed-state replica with antiperiodic Dirac gluing and the less singular extension. Let \(\Sigma\) be the fixed plane of a rotation under which \(\Theta\) is invariant, so that \(\Sigma\) is totally geodesic for \(G\). With unit normals \(n_i\) to \(\Sigma\), write \(R_{ii}=R_{\mu\nu}n^\mu_in^\nu_i\) and \(R_{ijij}=R_{\mu\nu\rho\sigma}n^\mu_in^\nu_jn^\rho_in^\sigma_j\) (sums over \(i,j\)), and \(R_\Sigma\) for the intrinsic scalar curvature of \(\Sigma\), so that \(R_\Sigma=R-2R_{ii}+R_{ijij}\) for totally geodesic \(\Sigma\). The pulled-back background is the same on every sheet, so \(\operatorname{Tr}\rho^n\) is the sum over the twist phases \(e^{2\pi\mathrm ik/n}\) of the single-sheet problem with a flat flavour vortex of flux \(a_k=k/n\) at \(\Sigma\) (integer \(k\) for the bosonic sectors, \(k=-\tfrac{n-1}2,\dots,\tfrac{n-1}2\) for the Dirac fields). For a Laplace-type sector \(-\nabla^2+E\) the \(s^0\) surface term is linear in the dimension-two data. For a component with a periodic phase it is the scalar term of Equation~(70) of~\cite{Solodukhin2011}, \(c_B(n)\operatorname{tr}(\tfrac R6-E)+d_B(n)(R_{ii}-2R_{ijij})\) with \(c_B=\tfrac\pi3\tfrac{1-n^2}n\) and \(d_B=-\tfrac\pi{180}\tfrac{(1-n^2)(1+n^2)}{n^3}\), and it has no term in the bundle curvature, because every periodic twist acts as a common phase and the sum of the Bernoulli polynomials \(B_3(k/n)\) in the linear coefficient vanishes. For a Dirac field the common phase is accompanied by the spin term \(\gamma^{12}\Omega_{12}\) of \((\gamma^\mu D_\mu)^2\), with \(\Omega_{12}=[D_1,D_2]\), whose part supported on \(\Sigma\) is attractive in the component where it is positive; the less singular extension then gives the \(E\)-linear weight \(-\tfrac12c_B\) per component and a term \(-\tfrac12\sum_ka_k^2\operatorname{tr}(\gamma^{12}\Omega_{12})\) with \(\sum_ka_k^2=(n^2-1)/(12n)\) and \(\Omega_{12}\) the bundle curvature in the normal plane (from the exact Landau--Aharonov--Bohm spectrum; the terms odd in \(a\) cancel in the symmetric sum over \(k\)). The remaining geometric part of the Dirac term is fixed at first order in \(n-1\) by the value \(a=11/360\) of the Dirac field (Equation~(275) of~\cite{Solodukhin2011}, after Duff; not re-derived here), which gives \(-\tfrac{11}{90}\) for a sphere. With one proper-time limit, the coefficient of \(\ln\LUV\) in the replica entropy is \(\int_\Sigma\sigma_{\log}\), where, with \(X'=-\operatorname{tr}(b_\sigma b^\sigma)\) in the vector sector, \(	au_\Omega=\operatorname{tr}(\gamma^{12}\Omega^{f}_{12})\) the normal-plane curvature of the Dirac connection, with the same \(\Omega_{12}=[D_1,D_2]\), \(e\) the scalar endomorphism of the preceding paragraph and \(W=W_{ijij}\) the Weyl projection,
\[
  24\pi\,\sigma_{\log}=\tfrac43R_G+2R_{ii}-2R_{ijij}+2W-12\,e+\tfrac32X'-\tau_\Omega .
\]
On the self-dual and anti-self-dual backgrounds, \(W=0\) and \(e=R_G/6\), so that \(24\pi\sigma_{\log}=\tfrac32X'-\tau_\Omega\). At three sampled points of \(\Sigma\) for the rotation-invariant self-dual example above this equals \(2R_\Sigma\), twice the intrinsic curvature of \(\Sigma\) (exact rational arithmetic; not proved), so that \(\sigma_{\log}=R_\Sigma/12\pi\). The cone logarithmic coefficient is then an Euler density of \(\Sigma\) with nonvanishing local value. The coefficient \(\propto\int\rho^2\mathcal S_4\) of the induced action vanishes on these backgrounds, and so does the first variation of \(\int\rho^2\mathcal S_4\) along the smooth replica, because \(\mathcal S_4\ge0\) vanishes there. The two sides therefore differ at logarithmic order, as coefficients, for the fixed-state replica on these backgrounds; this is a statement about the fixed-state replica only and not about \(G_{\mathrm{eff}}\). The first-order smooth replica of the preceding paragraphs gives zero on both sides, and which replica the matrix state selects is not derived here. Contact and edge-mode terms of the connection and Dirac sectors, which are not included, shift the coefficient of \(R_\Sigma\) and could cancel the difference.

These statements are for the Poisson-limit operators, the Euclidean section, one loop, \(\alpha=1\), \(N_s=6\), \(n_D=2\), the totally geodesic \(\Sigma\), the fixed-state replica and one common proper-time limit. They use Equation~(70) of~\cite{Solodukhin2011} and the Dirac coefficient \(a=11/360\) as inputs, do not treat finite terms, edge modes, other extensions, higher orders in \(n-1\), or the non-self-dual \(\mathfrak J=0\) backgrounds beyond the general formula, and give no value or sign of \(G_{\mathrm{eff}}\). Problem~B remains open and PROGRAMME-CONDITIONAL.
